\chapter{Building a Multi-Body Model}
\label{ch:building-models}

\epigraph{The model is the hypothesis made physical.}{}

\section{Model Formats}

\subsection{URDF (Universal Robot Description Format)}
XML-based format used by ROS, Drake, and Pinocchio. Defines links, joints, and visual/collision geometry.

\subsection{MJCF (MuJoCo XML)}
MuJoCo's native format. Extends URDF with actuators, tendons, muscles, and contact parameters.

\section{Building a 7-Segment Golf Model}

\begin{lstlisting}[language=Python, caption={Programmatic multi-body model construction.}]
import numpy as np
from dataclasses import dataclass, field

@dataclass
class Link:
    name: str
    mass: float
    length: float
    com_position: float
    inertia: np.ndarray = field(default_factory=lambda: np.eye(3))

@dataclass
class Joint:
    name: str
    parent: str
    child: str
    joint_type: str  # revolute, ball, fixed
    axis: list[float] = field(default_factory=lambda: [0, 0, 1])
    limits: tuple[float, float] = (-np.pi, np.pi)

@dataclass
class MultiBodyModel:
    name: str
    links: list[Link] = field(default_factory=list)
    joints: list[Joint] = field(default_factory=list)

    @property
    def n_dof(self) -> int:
        dof_map = {"revolute": 1, "ball": 3, "fixed": 0}
        return sum(dof_map.get(j.joint_type, 1) for j in self.joints)

    def to_urdf(self) -> str:
        """Export model as URDF XML string."""
        lines = [f'<robot name="{self.name}">']
        for link in self.links:
            lines.append(f'  <link name="{link.name}">')
            lines.append(f'    <inertial>')
            lines.append(f'      <mass value="{link.mass}"/>')
            lines.append(f'    </inertial>')
            lines.append(f'  </link>')
        for joint in self.joints:
            lines.append(f'  <joint name="{joint.name}" '
                         f'type="{joint.joint_type}">')
            lines.append(f'    <parent link="{joint.parent}"/>')
            lines.append(f'    <child link="{joint.child}"/>')
            lines.append(f'  </joint>')
        lines.append('</robot>')
        return '\n'.join(lines)


def build_golf_swing_model() -> MultiBodyModel:
    """Build a 7-segment golf swing model.

    Segments: pelvis, trunk, upper_arm, forearm, hand, club_shaft, club_head
    """
    model = MultiBodyModel(name="golf_swing_7seg")

    # Segment properties (de Leva 1996, 80kg male)
    model.links = [
        Link("pelvis", mass=11.2, length=0.15, com_position=0.075),
        Link("trunk", mass=25.5, length=0.48, com_position=0.22),
        Link("upper_arm", mass=2.16, length=0.28, com_position=0.16),
        Link("forearm", mass=1.30, length=0.27, com_position=0.12),
        Link("hand", mass=0.49, length=0.19, com_position=0.07),
        Link("club_shaft", mass=0.10, length=1.05, com_position=0.50),
        Link("club_head", mass=0.20, length=0.10, com_position=0.05),
    ]

    model.joints = [
        Joint("pelvis_rotation", "world", "pelvis", "revolute", [0, 1, 0]),
        Joint("trunk_flexion", "pelvis", "trunk", "revolute", [0, 0, 1]),
        Joint("shoulder", "trunk", "upper_arm", "ball"),
        Joint("elbow", "upper_arm", "forearm", "revolute", [0, 0, 1]),
        Joint("wrist", "forearm", "hand", "revolute", [0, 0, 1]),
        Joint("grip", "hand", "club_shaft", "fixed"),
    ]

    return model


model = build_golf_swing_model()
print(f"Golf swing model: {len(model.links)} segments, "
      f"{model.n_dof} DOF")
print(f"\nURDF preview:\n{model.to_urdf()[:500]}...")
\end{lstlisting}

\section*{Exercises}
\begin{enumerate}
  \item Build a 3-segment arm model and export as URDF.
  \item Add muscle attachment points to the golf model.
  \item Load the URDF into MuJoCo and verify the model visually.
\end{enumerate}
